\section{Methods: Rashomon Pipeline}

\subsection{Input}

The pipeline takes as input:
\begin{itemize}
    \item A dataset $(X, y)$ where $X$ is a feature matrix and $y$ is a binary target variable
    \item Feature metadata specifying which features are numeric vs.\ categorical
    \item A tolerance parameter $\varepsilon$ defining the Rashomon set threshold
    \item Model family configurations (Logistic Regression, Random Forest, Gradient Boosting, MLP)
    \item Hyperparameter search spaces for each model family
    \item Number of hyperparameter configurations to sample per family ($n_{\text{models}}$)
    \item Random seed(s) for reproducibility
\end{itemize}

\subsection{Model Training}

The pipeline trains multiple models across four model families:
\begin{itemize}
    \item \textbf{Logistic Regression (LogReg)}: L2 regularization, varying $C \in \{0.001, 0.01, 0.1, 1.0, 10.0, 100.0\}$, solvers: \texttt{lbfgs}, \texttt{liblinear}
    \item \textbf{Random Forest (RF)}: Varying \texttt{n\_estimators}, \texttt{max\_depth}, \texttt{min\_samples\_split}, \texttt{min\_samples\_leaf}, \texttt{max\_features}, \texttt{bootstrap}
    \item \textbf{Gradient Boosting (GBM)}: Varying \texttt{n\_estimators}, \texttt{learning\_rate}, \texttt{max\_depth}, \texttt{subsample}
    \item \textbf{MLP}: Varying \texttt{hidden\_layer\_sizes}, \texttt{alpha}, \texttt{learning\_rate\_init}, \texttt{activation}
\end{itemize}

For each model family, $n_{\text{models}}$ hyperparameter configurations are sampled uniformly at random from the search space. Each configuration is trained with one or more random seeds to account for stochasticity. Numeric features are standardized for LogReg and MLP; categorical features are one-hot encoded for all models.

\subsection{Definition of Near-Optimal (Rashomon Set)}

A model $m$ is considered \emph{near-optimal} and included in the Rashomon set if its validation loss $L(m)$ satisfies:
\begin{equation}
L(m) \leq L^* + \varepsilon
\end{equation}
where $L^* = \min_{m'} L(m')$ is the best validation loss across all trained models, and $\varepsilon$ is the tolerance parameter (typically $\varepsilon = 0.01$ for log loss, corresponding to approximately 1\% worse than the best model).

The Rashomon set can be constructed in two modes:
\begin{itemize}
    \item \textbf{Global}: Single threshold $L^* + \varepsilon$ applied across all model families
    \item \textbf{Per-family}: Separate threshold $L^*_f + \varepsilon$ for each model family $f$, where $L^*_f$ is the best loss within that family
\end{itemize}

\subsection{Cross-Validation Usage}

Cross-validation (CV) is used to compute model performance for Rashomon set membership, making the set less dependent on a single train/validation split. The pipeline supports two CV methods:

\begin{itemize}
    \item \textbf{K-fold CV}: $K$ folds (default $K=5$) on the training data
    \item \textbf{Repeated holdout}: $R$ random train/validation splits (default $R=5$)
\end{itemize}

When CV is enabled:
\begin{enumerate}
    \item The dataset is first split into a held-out test set (typically 20\%) and train+validation data
    \item CV splits are created on the train+validation data
    \item For each model, the mean CV loss $\bar{L}_{\text{CV}}(m) = \frac{1}{K} \sum_{k=1}^K L_k(m)$ is computed across folds
    \item Rashomon membership is determined using $\bar{L}_{\text{CV}}(m) \leq \bar{L}^*_{\text{CV}} + \varepsilon$
\end{enumerate}

This approach improves stability and reproducibility of the Rashomon set, as membership is no longer conditional on a single train/validation split.

\subsection{Test Set Usage}

The test set is held out from all model training and selection procedures. It is used exclusively for:
\begin{itemize}
    \item Generating final predictions $P_{\text{test}}$ from all models in the Rashomon set
    \item Computing observation-wise multiplicity metrics (variance, prediction intervals, flip instability)
    \item Spatial analysis (Moran's I, LISA hotspots)
    \item All downstream analyses of predictive multiplicity
\end{itemize}

The test set is \emph{not} used for:
\begin{itemize}
    \item Model training
    \item Hyperparameter selection
    \item Rashomon set membership decisions
\end{itemize}

\subsection{Pipeline Outputs}

The pipeline produces the following outputs:

\begin{enumerate}
    \item \textbf{Predictions matrix} $P_{\text{test}} \in \mathbb{R}^{M \times N}$: Predicted probabilities for all $M$ Rashomon models on all $N$ test observations. Saved as \texttt{P\_test.npy}.
    
    \item \textbf{Metadata} \texttt{meta.csv}: DataFrame containing for each Rashomon model:
    \begin{itemize}
        \item Model family name
        \item Hyperparameter configuration (as dictionary)
        \item Random seed used
        \item Validation loss (or mean CV loss if CV was used)
    \end{itemize}
    
    \item \textbf{Multiplicity metrics} \texttt{metrics.npz}: Observation-wise metrics computed from $P_{\text{test}}$:
    \begin{itemize}
        \item \texttt{variance}: $\text{Var}_m(\hat{p}_{mi})$ for each observation $i$
        \item \texttt{normalized\_variance}: Variance normalized by mean prediction
        \item \texttt{flip\_instability}: Fraction of models that flip the binary prediction at threshold 0.5
        \item \texttt{pi\_width\_90/95}: Width of 90\%/95\% prediction intervals
    \end{itemize}
    
    \item \textbf{Split information} \texttt{split.npz}: Indices for train, validation, and test sets (or test set + CV splits if CV was used)
    
    \item \textbf{Test data} \texttt{X\_test.csv}, \texttt{y\_test.npy}: Feature matrix and target labels for the test set
    
    \item \textbf{Configuration} \texttt{config.json}: All hyperparameters used in the experiment (dataset, $\varepsilon$, number of models, CV settings, etc.)
\end{enumerate}

These outputs enable downstream analyses including spatial hotspot detection, hyperparameter variance decomposition, and interpretable rule extraction.
